\chapter{Intersections, distances, and numerical tools}
\label{chap:algorithms}

An intersection query asks where a one-dimensional entity meets a circle,
sphere, or plane. A distance query asks for a closest point as well as its
parameter. These two questions share numerical machinery, but their results
have different meanings: an empty intersection set does not say how far apart
the entities are, and a distance rounded to zero does not create an intersection
record.

The examples in this chapter use \api{namespace np = nurbspath;} and
\api{using REAL = double;}, with the umbrella header included. All native 2D
queries stay entirely in the 2D world. An explicit \api{project(plane, entity)}
constructs a 3D entity when an application deliberately wants to cross that
boundary; it is not an implementation step of a 2D intersection or distance.

\section{Choose a query and its controls}

\begin{table}[htbp]
\centering
\small
\begin{tabularx}{\linewidth}{@{}l Y l@{}}
\toprule
World & Named intersection function & Result type \\
\midrule
2D & \api{intersect_ray_circle(ray, circle, settings)} & \api{intersection_result2<REAL>} \\
2D & \api{intersect_spline_circle(spline, circle, settings)} & \api{intersection_result2<REAL>} \\
3D & \api{intersect_ray_sphere(ray, sphere, settings)} & \api{intersection_result3<REAL>} \\
3D & \api{intersect_ray_plane(ray, plane, settings)} & \api{intersection_result3<REAL>} \\
3D & \api{intersect_spline_sphere(spline, sphere, settings)} & \api{intersection_result3<REAL>} \\
3D & \api{intersect_spline_plane(spline, plane, settings)} & \api{intersection_result3<REAL>} \\
\bottomrule
\end{tabularx}
\caption{The six supported intersection pairs. Each also has an
\api{intersect(first, second, settings)} convenience overload with the same
argument order and behavior.}
\label{tab:intersection-pairs}
\end{table}

The first argument is always the ray or ordinary NURBS spline. These overloads
do not provide arbitrary entity-pair intersections: there is no ray--ray,
spline--spline, sphere--sphere, circle--circle, or reversed-argument overload.
For an array-backed spline, first construct its corresponding ordinary spline
explicitly; that conversion validates the current public array storage.

\needspace{12\baselineskip}
\subsection{Two tolerances, two different units}

\api{numerical_settings<REAL>} carries all controls for numerical
intersections and point-to-spline distances. The defaults are:
\par\noindent\begin{minipage}{\linewidth}
\begin{lstlisting}[language=C++]
np::numerical_settings<REAL> settings;
// residual_tolerance = 1.0000e-08
// parameter_tolerance = 1.0000e-10
// max_iterations = 100
// sample_count = 256

settings.residual_tolerance = 1.0000e-07;
settings.parameter_tolerance = 1.0000e-09;
settings.sample_count = 512;
settings.validate();
\end{lstlisting}
\end{minipage}\par
\api{residual_tolerance} is measured in the units of the supplied function's
result. For the library's surface intersections, the residual is a signed
world distance; for \api{find_roots} it has whatever units the callable returns.
\api{parameter_tolerance} measures a parameter interval. For spline queries it
uses the original knot-domain units, not a normalized fraction of the domain.
Changing a spline's native domain therefore changes the interpretation of a
fixed parameter tolerance.

\api{sample_count} is the number of initial uniform intervals, so a sampled
search evaluates \api{sample_count + 1} initial positions. The iteration limit
applies to each refinement, rather than being a budget shared by all contacts
or local minima. Increase sampling density when a curve or residual oscillates
rapidly, has tightly spaced contacts, or has narrow distance minima.

\api{validate()} throws \api{std::invalid_argument} unless both tolerances are
finite and strictly positive, \api{max_iterations} is positive, and
\api{sample_count} is at least two. The numerical query functions call this
validation themselves. A larger iteration budget sharpens a basin that was
found; it cannot recover a basin that the initial sampling missed.

\begin{manualnote}[What a sampled answer guarantees]
Sampling and local refinement provide useful coverage controlled by
\api{sample_count}. They do not establish an unconditional continuous global
guarantee for every NURBS curve or user-supplied function. Check stability under
increased sampling when the application requires confidence that all contacts
or the smallest distance have been found.
\end{manualnote}

\section{Read an intersection result}

\api{intersection_kind} has three values. \api{none} means that the query found
no contact in its searched domain; \api{discrete} means that isolated contacts
are listed; \api{coincident} denotes a continuous overlap that a finite list of
points cannot represent. A coincident result deliberately has an empty
\api{points} vector. Thus \api{points.empty()} alone cannot distinguish no
contact from continuous overlap. Use \api{has_intersection()} to accept either
kind of contact, and \api{is_coincident()} to branch on continuous overlap.

\begin{table}[htbp]
\centering
\small
\begin{tabularx}{\linewidth}{@{}l Y Y@{}}
\toprule
Field & \api{intersection2<REAL>} & \api{intersection3<REAL>} \\
\midrule
\api{point} & Contact position in 2D & Contact position in 3D \\
\api{s} & Native ray or spline parameter & Native ray or spline parameter \\
\api{u} & Circle angle in $[0,2\pi)$ & Sphere longitude or plane $u$ coordinate \\
\api{v} & Not present & Sphere latitude or plane $v$ coordinate \\
\api{residual} & Absolute circle-distance residual & Absolute sphere- or plane-distance residual \\
\bottomrule
\end{tabularx}
\caption{A contact preserves both the world position and the participating
entities' parameters. Sphere latitude lies in $[-\pi/2,\pi/2]$; plane parameters
are lengths along its orthonormal local frame.}
\end{table}

The vector \api{points} is sorted by increasing \api{s}. Numerically close
root parameters are merged; the root solver's merge threshold is the greater
of four parameter tolerances and a roundoff allowance proportional to the
search interval. This is a parameter-based rule, not a promise to merge every
pair of spatially close points. A self-intersecting spline can visit the same
world position at two distinct native parameters, and those visits can remain
two distinct contacts.

\begin{figure}[htbp]
\centering
\begin{minipage}{0.48\linewidth}
\centering
\begin{tikzpicture}
\begin{axis}[width=\linewidth,height=5.0cm,axis equal image,
 xmin=-3.3,xmax=1.5,ymin=-1.45,ymax=1.55,
 axis lines=middle,xlabel={$x$},ylabel={$y$},
 ticks=none,clip=false]
\addplot[manualblue,thick,domain=0:360,samples=97]
 ({cos(x)},{sin(x)});
\addplot[manualred,thick,-{Stealth}] coordinates {(-3,0) (1.35,0)};
\addplot[only marks,mark=*,mark size=2.4pt,manualink]
 table[x index=1,y index=2] {pictures/intersection_crossing.dat};
\node[below left] at (axis cs:-3,0) {$O$};
\node[below] at (axis cs:-1,0) {$s=1$};
\node[below] at (axis cs:1,0) {$s=2$};
\node[above] at (axis cs:-1.9,0) {$\mathbf d=(2,0)$};
\end{axis}
\end{tikzpicture}
\end{minipage}\hfill
\begin{minipage}{0.48\linewidth}
\centering
\begin{tikzpicture}
\begin{axis}[width=\linewidth,height=5.0cm,axis equal image,
 xmin=-2.4,xmax=1.5,ymin=-1.45,ymax=1.55,
 axis lines=middle,xlabel={$x$},ylabel={$y$},
 ticks=none,clip=false]
\addplot[manualblue,thick,domain=0:360,samples=97]
 ({cos(x)},{sin(x)});
\addplot[manualred,thick,-{Stealth}] coordinates {(-2,1) (1.35,1)};
\addplot[only marks,mark=*,mark size=2.4pt,manualink]
 table[x index=1,y index=2] {pictures/intersection_tangent.dat};
\node[above left] at (axis cs:-2,1) {$O$};
\node[above right] at (axis cs:0,1) {$s=2$};
\node[below left] at (axis cs:-0.75,1) {$\mathbf d=(1,0)$};
\end{axis}
\end{tikzpicture}
\end{minipage}\par
\caption{Crossing and tangent ray--circle contacts. The black points are
computed by the library; TikZ constructs the unit circles and rays. The tangent
contact has no residual sign change.}
\label{fig:ray-circle-contacts}
\end{figure}

\section{Rays: forward domains and direction scales}

\subsection{A ray through a circle or sphere}

The ray equation is $R(s)=O+s\mathbf d$, with $s\geq0$. A direction is
not normalized by the entity. For the left-hand example in
Figure~\ref{fig:ray-circle-contacts}, each increase of one in $s$ moves two
world units. This gives the two unit-circle contacts at $s=1$ and $s=2$.
\par\noindent\begin{minipage}{\linewidth}
\begin{lstlisting}[language=C++]
const np::ray2<REAL> ray{{-3, 0}, {2, 0}};
const np::circle2<REAL> circle{{0, 0}, 1};
const auto hits = np::intersect_ray_circle(ray, circle);
// Equivalent: np::intersect(ray, circle)

std::cout << std::scientific << std::setprecision(4);
if (hits.kind == np::intersection_kind::discrete) {
    for (const auto& hit : hits.points) {
        std::cout << hit.s << ' ' << hit.point.x << ' '
                  << hit.point.y << '\n';
    }
}
\end{lstlisting}
\end{minipage}\par
The rounded output is:
\par\noindent\begin{minipage}{\linewidth}
\begin{lstlisting}
1.0000e+00 -1.0000e+00 0.0000e+00
2.0000e+00  1.0000e+00 0.0000e+00
\end{lstlisting}
\end{minipage}\par
The 3D counterpart changes the target from a circle to a sphere:
\par\noindent\begin{minipage}{\linewidth}
\begin{lstlisting}[language=C++]
const np::ray3<REAL> ray3{{-3, 0, 0}, {2, 0, 0}};
const np::sphere3<REAL> sphere{{0, 0, 0}, 1};
const auto sphere_hits = np::intersect_ray_sphere(ray3, sphere);
// Equivalent: np::intersect(ray3, sphere)
\end{lstlisting}
\end{minipage}\par
Both radial queries refine the signed distance
\[
 f(t)=\|O+t\widehat{\mathbf d}-C\|-r
\]
using world-distance travel $t$, then return the native parameter
$s=t/\|\mathbf d\|$. The finite forward search extends to
$\|O-C\|+r+\varepsilon_r$. Scaling the direction by a positive factor therefore
preserves world contacts while dividing their native $s$ values by that factor.
A negative factor reverses the forward half-line and can change the set of
contacts. A ray starting inside a circle or sphere normally has one forward
exit; a ray pointing away from the target may have none. Contact at the ray
origin is included.

\subsection{A ray and an infinite plane}

\api{intersect_ray_plane} uses the plane's unit normal and the signed
world-distance residual $f(s)=\widehat{\mathbf n}\mathbin{\cdot}(R(s)-P_0)$.
The native rate of change is $\mathbf d\mathbin{\cdot}\widehat{\mathbf n}$.
If that rate is exactly zero, an origin within residual tolerance gives
\api{coincident}; otherwise the result is \api{none}. This parallel test uses
a zero normal component rather than comparing the direction's magnitude to a
tolerance.

For a nonzero rate, an origin already within tolerance is an isolated contact
at $s=0$. Otherwise a negative estimated contact parameter means that the
plane is behind the ray. A forward contact is bracketed and refined through
the shared numerical root search. Refinement tightens the parameter tolerance
where needed for a large normal-direction rate.
\par\noindent\begin{minipage}{\linewidth}
\begin{lstlisting}[language=C++]
const np::plane3<REAL> plane{{0, 0, 0}, {0, 0, 1}};
const np::ray3<REAL> toward{{0.5, 0.25, 2}, {0, 0, -2}};
const auto plane_hit = np::intersect_ray_plane(toward, plane);
// Equivalent: np::intersect(toward, plane)
// One contact: s = 1, world point (0.5, 0.25, 0).

const np::ray3<REAL> parallel{{0, 0, 2}, {1, 0, 0}};
const np::ray3<REAL> contained{{0, 0, 0}, {1, 0, 0}};
const auto no_hit = np::intersect(parallel, plane);
const auto overlap = np::intersect(contained, plane);
// no_hit.kind == intersection_kind::none
// overlap.is_coincident() && overlap.points.empty()
\end{lstlisting}
\end{minipage}\par

\begin{figure}[htbp]
\centering
\begin{tikzpicture}[x=1cm,y=1cm,font=\small]
\foreach \x/\title in {0/{forward crossing},4.7/{parallel, separated},9.4/{contained ray}} {
  \begin{scope}[shift={(\x,0)}]
    \fill[manualteal!10] (-1.65,-0.18) -- (0.95,-0.18) -- (1.65,0.35) -- (-0.95,0.35) -- cycle;
    \draw[manualteal] (-1.65,-0.18) -- (0.95,-0.18) -- (1.65,0.35) -- (-0.95,0.35) -- cycle;
    \node at (0,-0.72) {\title};
  \end{scope}
}
\draw[manualred,thick,-{Stealth}] (-0.7,1.65) -- (0.55,-0.7);
\fill[manualink] (0.18,0) circle (2pt);
\draw[manualred,thick,-{Stealth}] (3.25,1.2) -- (6.15,1.2);
\draw[manualred,thick,-{Stealth}] (8.0,0.08) -- (10.65,0.08);
\node[manualmuted] at (0,2) {\api{discrete}};
\node[manualmuted] at (4.7,2) {\api{none}};
\node[manualmuted] at (9.4,2) {\api{coincident}};
\end{tikzpicture}
\caption{The three ray--plane classifications. The drawn patch represents an
infinite plane; its finite outline is only a diagram aid.}
\label{fig:ray-plane-cases}
\end{figure}

\section{Splines: crossings, tangencies, and continuous contact}

The three spline queries evaluate the full active interval
$[s_{\min},s_{\max}]$. Their residuals follow the target surface:
\begin{itemize}
\item \api{intersect_spline_circle} uses a 2D radial distance.
\item \api{intersect_spline_sphere} uses a 3D radial distance.
\item \api{intersect_spline_plane} uses the plane's signed distance.
\end{itemize}
Their \api{intersect} aliases preserve exactly these domains and meanings.

A compact degree-one example makes the native domain visible:
\par\noindent\begin{minipage}{\linewidth}
\begin{lstlisting}[language=C++]
const np::nurbs_spline2<REAL> segment2{
    std::vector<np::point2<REAL>>{{-2, 0}, {2, 0}},
    std::vector<REAL>{1, 1},
    std::vector<REAL>{10, 10, 14, 14}, 1};
const np::circle2<REAL> unit_circle{{0, 0}, 1};
const auto circle_hits = np::intersect_spline_circle(
    segment2, unit_circle, settings);
// Native s values are approximately 11 and 13, not 0.25 and 0.75.

const np::nurbs_spline3<REAL> segment3{
    std::vector<np::point3<REAL>>{{-2, 0, 0}, {2, 0, 0}},
    std::vector<REAL>{1, 1},
    std::vector<REAL>{10, 10, 14, 14}, 1};
const np::sphere3<REAL> unit_sphere{{0, 0, 0}, 1};
const np::plane3<REAL> yz_plane{{0, 0, 0}, {1, 0, 0}};
const auto sphere_contacts = np::intersect_spline_sphere(
    segment3, unit_sphere, settings);
const auto plane_contact = np::intersect_spline_plane(
    segment3, yz_plane, settings);
// Sphere: s approximately 11 and 13. Plane: s approximately 12.
\end{lstlisting}
\end{minipage}\par
Endpoints are included. A tangent contact may have residual zero with the same
sign on both sides, so the search does more than look for sign changes. It also
refines sampled local minima of the absolute residual. Increasing
\api{sample_count} remains necessary when tangencies lie in narrow or crowded
features.

Before returning isolated roots, each spline query checks whether the surface
residual stays within tolerance at a uniform set of
$\max(\api{sample_count},16)+1$ points. If so, it returns \api{coincident}.
This classification is numerical and sampled: it is not an algebraic proof
that the entire rational curve lies on the surface. A spline contained in a
plane or tracing a circle on a sphere is a natural use case.

For a spline explicitly marked closed, a root at the last active parameter is
removed when roots are already present within parameter tolerance of both
domain ends. The first-end contact remains as the seam representative. Closure
does not mean that arbitrary equal world positions elsewhere on the curve are
merged, and it does not itself imply periodic derivatives.

\begin{figure}[htbp]
\centering
\begin{tikzpicture}[x=1cm,y=1cm,font=\small]
\begin{scope}
\draw[manualblue,thick] (0,0) circle (1.35);
\draw[manualteal,thick] (-1.85,0) -- (1.85,0);
\fill[manualink] (1.35,0) circle (2.6pt);
\fill[manualink] (-1.35,0) circle (2.6pt);
\draw[manualorange,-{Stealth}] (1.05,0.85) arc[start angle=39,end angle=115,radius=1.35];
\node[anchor=west,align=left] at (1.55,0.7) {seam: same position\\$s_{\min}$ retained\\$s_{\max}$ suppressed};
\draw[manualmuted] (1.45,0.48) -- (1.35,0.08);
\node at (0,-1.85) {closed curve meeting a surface};
\end{scope}
\begin{scope}[shift={(8.6,0)}]
\fill[manualteal!10] (-1.8,-0.9) -- (1.2,-0.9) -- (1.9,0.8) -- (-1.1,0.8) -- cycle;
\draw[manualteal] (-1.8,-0.9) -- (1.2,-0.9) -- (1.9,0.8) -- (-1.1,0.8) -- cycle;
\draw[manualblue,very thick] (-1.25,-0.4) .. controls (-0.6,0.8) and (0.1,-0.75) .. (1.25,0.2);
\foreach \x/\y in {-1.25/-0.4,-0.75/0.09,-0.2/0.03,0.35/-0.03,0.9/0.07,1.25/0.2}
 \fill[manualink] (\x,\y) circle (1.6pt);
\node[align=center] at (0,1.35) {sampled residuals near zero\\\api{coincident}, no finite point list};
\node at (0,-1.85) {curve contained in a plane};
\end{scope}
\end{tikzpicture}
\caption{Closed-seam deduplication and sampled continuous-contact detection.
The left-hand surface is shown by its trace through the curve's drawing
plane. The right-hand dots illustrate residual checks rather than isolated
intersection records.}
\label{fig:seam-coincidence}
\end{figure}

\section{Distances return geometry, not just a number}

The detailed queries in \api{distance.hpp} keep the closest world point and
its target parameter. The corresponding \api{distance} overload calls the
detailed function and returns only its \api{distance} field. Argument order is
always query point first, target entity second.

\begin{table}[htbp]
\centering
\small
\begin{tabularx}{\linewidth}{@{}l Y Y@{}}
\toprule
World & Detailed query & Result and target parameters \\
\midrule
2D & \api{distance_to_circle(point, circle, tolerance)} & \api{point_circle_distance2<REAL>}: \api{u} \\
2D & \api{distance_to_spline(point, spline, settings)} & \api{point_spline_distance2<REAL>}: \api{s} \\
3D & \api{distance_to_plane(point, plane, tolerance)} & \api{point_surface_distance3<REAL>}: \api{u}, \api{v} \\
3D & \api{distance_to_sphere(point, sphere, tolerance)} & \api{point_surface_distance3<REAL>}: \api{u}, \api{v} \\
3D & \api{distance_to_spline(point, spline, settings)} & \api{point_spline_distance3<REAL>}: \api{s} \\
\bottomrule
\end{tabularx}
\caption{Every result also has \api{distance} and \api{closest_point} fields.
The scalar alternative for each row is
\api{distance(point, target, tolerance_or_settings)}.}
\label{tab:distance-pairs}
\end{table}

All reported distances are nonnegative. When the measured separation is at
most the contact tolerance, the returned scalar is snapped to zero; the closest
point and its parameters remain those calculated on the target. A zero result
therefore permits a small geometric gap. For a spline, the contact tolerance
is \api{settings.residual_tolerance}; for a circle, sphere, or plane, the
separate tolerance argument defaults to \api{1.0000e-08}.

\subsection{Circle and sphere: radial closest points}

For a center $C$, radius $r$, and query $Q\ne C$,
\[
 P=C+r\frac{Q-C}{\|Q-C\|},\qquad
 d_{\mathrm{measured}}=\big|\|Q-C\|-r\big|.
\]
The expression applies in 2D to a circle and in 3D to a sphere. It finds the
closest point on the boundary, including when $Q$ is inside. At the center,
every boundary point is equally near; the library deterministically chooses
the positive world-$x$ point. Its circle angle is $u=0$, and its spherical
parameters are $(u,v)=(0,0)$.
\par\noindent\begin{minipage}{\linewidth}
\begin{lstlisting}[language=C++]
const np::circle2<REAL> circle{{0, 0}, 1};
const auto radial2 = np::distance_to_circle(
    np::point2<REAL>{3, 4}, circle);
// Rounded closest_point = (6.0000e-01, 8.0000e-01).
// Rounded distance = 4.0000e+00.
const REAL scalar2 = np::distance(np::point2<REAL>{3, 4}, circle);

const np::sphere3<REAL> sphere{{0, 0, 0}, 1};
const auto radial3 = np::distance_to_sphere(
    np::point3<REAL>{0, 0, 3}, sphere);
// closest_point = (0, 0, 1), distance = 2, v = pi/2.
const REAL scalar3 = np::distance(np::point3<REAL>{0, 0, 3}, sphere);
\end{lstlisting}
\end{minipage}\par
The circle-distance tolerance must be finite and strictly positive. The
plane- and sphere-distance overloads reject nonpositive tolerances; applications
should also supply finite values, consistently with the settings-based queries.

\subsection{Plane: orthogonal closest points}

Let $\widehat{\mathbf n}$ be the unit normal and $P_0$ its parameter origin.
Then
\[
 h=\widehat{\mathbf n}\mathbin{\cdot}(Q-P_0),\qquad
 P=Q-h\widehat{\mathbf n},\qquad d_{\mathrm{measured}}=|h|.
\]
The detailed query returns plane parameters of $P$, so
\api{plane.point_at(result.u, result.v)} reconstructs the closest position.
\api{plane.signed_distance_to(point)} remains available when the sign, rather
than the nonnegative distance query, is required.
\par\noindent\begin{minipage}{\linewidth}
\begin{lstlisting}[language=C++]
const np::plane3<REAL> plane{
    np::point3<REAL>{0, 0, 0}, np::vector3<REAL>{0, 0, 1},
    np::vector3<REAL>{1, 0, 0}};
const np::point3<REAL> query{2, 3, 4};
const auto nearest = np::distance_to_plane(query, plane);
// closest_point = (2, 3, 0), distance = 4, (u,v) = (2,3).
const REAL separation = np::distance(query, plane);
\end{lstlisting}
\end{minipage}\par

\begin{figure}[htbp]
\centering
\begin{tikzpicture}[x=1cm,y=1cm,font=\small]
\begin{scope}
\draw[manualblue,thick] (0,0) circle (1.1);
\fill[manualink] (0,0) circle (2pt) node[below left] {$C$};
\draw[manualmuted,densely dotted] (0,0) -- (0.66,0.88);
\draw[manualred,thick] (0.66,0.88) -- (1.65,2.2);
\fill[manualink] (0.66,0.88) circle (2pt) node[below right] {$P$};
\fill[manualred] (1.65,2.2) circle (2pt) node[above right] {$Q$};
\node[above left] at (1.15,1.55) {$d$};
\fill[manualorange] (-0.33,0.44) circle (2pt) node[left] {$Q_{\rm inside}$};
\draw[manualorange,thick] (-0.33,0.44) -- (-0.66,0.88);
\fill[manualink] (-0.66,0.88) circle (2pt);
\node[align=center] at (0,-1.8) {radial distance to the boundary\\circle or sphere cross-section};
\end{scope}
\begin{scope}[shift={(8,0)}]
\fill[manualteal!10] (-1.9,-0.65) -- (1,-0.65) -- (2,0.5) -- (-0.9,0.5) -- cycle;
\draw[manualteal] (-1.9,-0.65) -- (1,-0.65) -- (2,0.5) -- (-0.9,0.5) -- cycle;
\draw[manualred,thick] (0,0) -- (0,2.1);
\draw[manualmuted] (0.2,0) -- (0.2,0.2) -- (0,0.2);
\fill[manualink] (0,0) circle (2pt) node[below] {$P$};
\fill[manualred] (0,2.1) circle (2pt) node[above] {$Q$};
\draw[manualteal,-{Stealth}] (1.2,0.1) -- (1.2,1.2) node[above] {$\widehat{\mathbf n}$};
\node[right] at (0.05,1.1) {$|h|$};
\node[align=center] at (0,-1.8) {orthogonal distance to an infinite plane\\$(u,v)$ describe the foot point $P$};
\end{scope}
\end{tikzpicture}
\caption{Analytic closest points for radial surfaces and planes. These
distances need no sampling. Interior radial queries still measure to the
surface, rather than returning zero throughout the filled region.}
\label{fig:analytic-distances}
\end{figure}

\section{Closest points on NURBS splines}

\api{distance_to_spline} minimizes the squared Euclidean separation
$D(s)=\|Q-C(s)\|^2$ over the active knot domain. Both 2D and 3D forms evaluate
the initial samples, retain the best sample, then refine each sampled local
basin using golden-section minimization. The first and last sample are also
eligible basins, with refinement over the adjacent boundary interval. This
matters when the true nearest point lies just inside an endpoint.

The same rational curve used in Chapter~\ref{chap:splines} provides a more
interesting distance example than a straight line:
\par\noindent\begin{minipage}{\linewidth}
\begin{lstlisting}[language=C++]
const np::nurbs_spline2<REAL> curve{
    std::vector<np::point2<REAL>>{
        {0, 0}, {1, 2}, {2, -1}, {3, 2}, {4, 0}},
    std::vector<REAL>{1, 1, 0.7, 1.5, 1},
    std::vector<REAL>{0, 0, 0, 0, 2, 5, 5, 5, 5}, 3};
const np::point2<REAL> query{2, 2.5};
const auto closest = np::distance_to_spline(query, curve, settings);
const REAL only_distance = np::distance(query, curve, settings);

// closest.s stays in the original [0,5] knot domain.
// curve.evaluate(closest.s) reproduces closest.closest_point.
\end{lstlisting}
\end{minipage}\par
Figure~\ref{fig:spline-closest} uses points evaluated by the library itself.
The target parameter remains the native $s$, even if its units were inherited
from interpolation stations. Neither the reported closest point nor the
distance query renormalizes the knot vector.

\begin{figure}[htbp]
\centering
\begin{tikzpicture}
\begin{axis}[width=0.88\linewidth,height=6.0cm,axis equal image,
 xmin=-0.2,xmax=4.2,ymin=-0.3,ymax=2.9,
 axis lines=left,xlabel={$x$},ylabel={$y$},
 xtick={0,1,2,3,4},ytick={0,1,2},grid=major,
 grid style={manualmuted!15},legend style={draw=none,font=\small},
 legend pos=south east]
\addplot[manualblue,very thick] table[x index=1,y index=2] {pictures/spline_curve.dat};
\addlegendentry{rational spline $C(s)$}
\addplot[manualred,thick,densely dashed,mark=*,mark size=2.3pt]
 table[x index=0,y index=1] {pictures/distance_reference.dat};
\addlegendentry{query-to-closest segment}
\node[above] at (axis cs:2,2.5) {$Q=(2,2.5)$};
\end{axis}
\end{tikzpicture}
\caption{A sampled and refined 2D closest point. The curve samples and the two
endpoints of the red distance segment come from the library reference-data
generator in \api{manual/pictures/}. The 3D distance overload follows the same
search strategy in its independent world.}
\label{fig:spline-closest}
\end{figure}

The 3D call has the same shape:
\par\noindent\begin{minipage}{\linewidth}
\begin{lstlisting}[language=C++]
const np::nurbs_spline3<REAL> segment{
    std::vector<np::point3<REAL>>{{0, 0, 0}, {1, 0, 0}},
    std::vector<REAL>{1, 1},
    std::vector<REAL>{0, 0, 1, 1}, 1};
const auto near_start = np::distance_to_spline(
    np::point3<REAL>{0.001, 0, 0}, segment, settings);
// closest point is just inside the segment, not its start endpoint.
// Reported distance is zero within residual_tolerance.
\end{lstlisting}
\end{minipage}\par
Endpoint retention prevents a refinement from discarding a true boundary
minimum, while boundary-interval refinement can improve a coarse endpoint
sample to an interior minimum. Multiple sampled minima are compared; no unique
closest-point guarantee is made for equal minima. In particular, a closed seam
or a symmetric curve can have more than one equally good parameter, while the
result type intentionally returns one.

\begin{figure}[htbp]
\centering
\begin{tikzpicture}[x=1cm,y=1cm,font=\small]
\draw[manualmuted,-{Stealth}] (-0.2,0) -- (11.8,0) node[right] {$s$};
\draw[manualmuted,-{Stealth}] (0,-0.1) -- (0,2.35) node[above] {$D(s)$};
\draw[manualblue,very thick] plot[smooth] coordinates {
 (0,0.75) (0.35,0.30) (0.7,0.22) (1.2,0.48) (2,1.4)
 (3,1.95) (4,1.15) (5,0.8) (6,1.05) (7,1.8)
 (8,1.5) (9,0.60) (10,0.30) (11,0.55)};
\foreach \x/\y in {0/0.75,2/1.4,4/1.15,6/1.05,8/1.5,10/0.30,11/0.55}
 \fill[manualink] (\x,\y) circle (2pt);
\draw[manualorange,line width=2pt] (0,-0.24) -- (2,-0.24);
\draw[manualteal,line width=2pt] (8,-0.24) -- (11,-0.24);
\fill[manualred] (0.63,0.21) circle (2.5pt);
\draw[manualred,-{Stealth}] (1.2,2.25) -- (0.65,0.35);
\node[anchor=west] at (1.2,2.25) {refine boundary interval};
\node[below,align=center] at (1,-0.32) {first sample can be a local basin\\despite an interior true minimum};
\node[below,align=center] at (9.5,-0.32) {refine another sampled basin\\compare candidate distances};
\end{tikzpicture}
\caption{Why closest-point searches refine more than the single best sample.
This schematic squared-distance graph illustrates a minimum near the beginning
of the domain and another basin near its end. It is not an additional library
test curve.}
\label{fig:distance-basins}
\end{figure}

\section{Reusable scalar algorithms}

The utilities in \api{utility.hpp} are public and independent of a particular
geometry class. Their callable parameters accept a scalar and return a scalar.
The application is responsible for choosing meaningful units, finite working
ranges, and functions compatible with the search assumptions.

\subsection{Find all sampled roots, including even-multiplicity contacts}

\api{find_roots<REAL>(function, lower, upper, settings)} requires a strictly
increasing interval. It samples that interval, refines sign changes with
bisection, and minimizes absolute residual near sampled sign-preserving local
minima. Contiguous within-tolerance sample bands are represented by a single
best candidate, avoiding a separate result for every sample on a flat contact.
The final roots are sorted and tolerance-deduplicated.
\par\noindent\begin{minipage}{\linewidth}
\begin{lstlisting}[language=C++]
const auto function = [](REAL x) {
    return (x + 0.8) * np::square(x - 0.3) * (x - 1.2);
};
const auto roots = np::find_roots<REAL>(
    function, -1.5, 1.6, settings);
for (REAL root : roots) {
    std::cout << std::scientific << std::setprecision(4)
              << root << '\n';
}
\end{lstlisting}
\end{minipage}\par
The three expected roots, shown as rounded reference values, are
$-8.0000\times10^{-1}$, $3.0000\times10^{-1}$, and
$1.2000\times10^{0}$. The middle root is double and does not change the sign
of the residual. \api{find_roots} rejects invalid settings or a nonpositive
interval length with \api{std::invalid_argument}; a nonfinite initial sample
value raises \api{std::domain_error}. A discontinuous callable or a nonfinite
value encountered during refinement falls outside a well-behaved root-search
problem; do not treat a sign change across a discontinuity as a root proof.

\begin{figure}[htbp]
\centering
\begin{tikzpicture}
\begin{axis}[width=0.88\linewidth,height=5.7cm,
 xmin=-1.5,xmax=1.6,ymin=-0.7,ymax=2.1,
 axis lines=middle,xlabel={$x$},ylabel={$f(x)$},
 xtick={-0.8,0.3,1.2},xticklabels={$-0.8$,$0.3$,$1.2$},
 ytick={0,1,2},grid=major,grid style={manualmuted!15},clip=true]
\addplot[manualblue,very thick] table[x index=0,y index=1] {pictures/algorithmroots.dat};
\addplot[only marks,mark=*,mark size=2.6pt,manualred]
 coordinates {(-0.8,0) (0.3,0) (1.2,0)};
\node[anchor=south,align=center] at (axis cs:0.3,0.2) {double root\\no sign change};
\draw[manualred,-{Stealth}] (axis cs:0.3,0.18) -- (axis cs:0.3,0.035);
\end{axis}
\end{tikzpicture}
\caption{The scalar residual
$f(x)=(x+0.8)(x-0.3)^2(x-1.2)$. Its sampled graph is generated by C++ using
the library's scalar helper; TikZ marks the known reference roots. Sign changes
alone would miss the contact at $x=0.3$.}
\label{fig:scalar-roots}
\end{figure}

\subsection{Refine one sign-changing bracket}

\api{bisect_root<REAL>(function, lower, upper, settings)} accepts endpoints in
either order. An endpoint whose absolute residual is already within tolerance
is returned immediately. Otherwise the endpoints must have opposite residual
signs, or \api{std::invalid_argument} is thrown. Each iteration keeps the
half-interval that still brackets the sign change. The function returns when
the residual is small enough, the interval is narrow enough, or the iteration
budget has been used.
\par\noindent\begin{minipage}{\linewidth}
\begin{lstlisting}[language=C++]
const REAL root_two = np::bisect_root<REAL>(
    [](REAL x) { return np::square(x) - 2; }, 1, 2, settings);
// Rounded: 1.4142e+00
\end{lstlisting}
\end{minipage}\par
The return type is a parameter alone: it does not carry an iteration count or
a convergence flag. In particular, reaching the parameter tolerance or the
iteration limit does not independently establish a small residual. Evaluate
the returned residual if the application requires that additional check.
Bisection is appropriate for a continuous sign-changing bracket; use the
sampled root finder to seek an unbracketed tangent contact.

\subsection{Minimize one bounded basin}

\api{golden_section_minimize<REAL>(function, lower, upper, settings)} also
accepts reversed endpoints. Two interior probes progressively contract a
closed interval, reusing one function value after each comparison. It returns
the center of the final interval, stopping at the parameter tolerance or
iteration limit. A unimodal objective over the chosen interval is the useful
working assumption. The function does not search every local minimum of an
arbitrary multi-basin objective, and it does not separately select endpoint
values; callers needing boundary minima should compare the endpoints too.
\par\noindent\begin{minipage}{\linewidth}
\begin{lstlisting}[language=C++]
const REAL minimum_x = np::golden_section_minimize<REAL>(
    [](REAL x) { return np::square(x - 0.75) + 2; },
    -1, 2, settings);
// Rounded minimizer: 7.5000e-01; objective value: 2.0000e+00.
\end{lstlisting}
\end{minipage}\par
Both refiners call \api{settings.validate()}. Golden-section search uses
\api{parameter_tolerance} and \api{max_iterations} to stop; changing
\api{residual_tolerance} does not supply a separate objective-value stopping
rule. It is nevertheless validated as part of the settings object.

\begin{figure}[htbp]
\centering
\begin{tikzpicture}[x=1cm,y=1cm,font=\small]
\begin{scope}
\draw[manualmuted,-{Stealth}] (-0.3,0) -- (5.8,0);
\draw[manualblue,thick] (-0.05,-0.8) .. controls (1.3,-0.7) and (3.8,0.9) .. (5.45,1.1);
\draw[manualink,thick] (0,-0.12) -- (0,0.12) node[above=3pt] {$a$};
\draw[manualink,thick] (5.4,-0.12) -- (5.4,0.12) node[below=3pt] {$b$};
\draw[manualorange,densely dashed] (2.7,-0.95) -- (2.7,1.05);
\node[below] at (2.7,-0.95) {$m=(a+b)/2$};
\draw[manualteal,line width=2pt] (0,-1.65) -- (2.7,-1.65);
\node[below,align=center] at (2.7,-1.75) {retain the sign-changing half\\bisection};
\end{scope}
\begin{scope}[shift={(7,0)}]
\draw[manualmuted,-{Stealth}] (-0.3,0) -- (5.8,0);
\draw[manualblue,thick] (0,1.3) .. controls (1.2,-0.8) and (2.7,-1.0) .. (5.4,1.25);
\draw[manualink,thick] (0,-0.12) -- (0,0.12) node[below=3pt] {$a$};
\draw[manualink,thick] (5.4,-0.12) -- (5.4,0.12) node[below=3pt] {$b$};
\draw[manualorange,densely dashed] (2.06,-1) -- (2.06,0.65);
\draw[manualorange,densely dashed] (3.34,-1) -- (3.34,0.65);
\node[above] at (2.06,0.65) {$\ell$};
\node[above] at (3.34,0.65) {$r$};
\draw[manualteal,line width=2pt] (0,-1.65) -- (3.34,-1.65);
\node[below,align=center] at (2.7,-1.75) {if $f(\ell)\leq f(r)$, retain $[a,r]$\\golden-section minimization};
\end{scope}
\end{tikzpicture}
\caption{A root refiner and a minimum refiner answer different questions.
The green segments show one possible retained interval. The diagrams are
schematic; neither refiner reports a continuous global-search certificate.}
\label{fig:scalar-refiners}
\end{figure}

\section{Linear systems and small helpers}

\subsection{Solve a square system}

\api{solve_linear_system(matrix, right_hand_side, pivot_tolerance)} solves
$A\mathbf x=\mathbf b$ by elimination with scaled partial pivoting. Each row's
largest absolute entry supplies a scale, so the pivot comparison is relative
to that row's magnitude. The matrix and right-hand-side arguments are passed
by value and copied for elimination; the caller's containers are preserved.

\par\noindent\begin{minipage}{\linewidth}
\begin{lstlisting}[language=C++]
const std::vector<std::vector<REAL>> matrix{{2, 1}, {1, 3}};
const std::vector<REAL> right_hand_side{5, 10};
const auto solution = np::solve_linear_system(matrix, right_hand_side);
// Rounded solution: {1.0000e+00, 3.0000e+00}.
\end{lstlisting}
\end{minipage}\par
The default pivot rejection threshold is
$64\,\api{std::numeric_limits<REAL>::epsilon()}$. The function requires a
nonempty square matrix and a matching right-hand side, otherwise it throws
\api{std::invalid_argument}. A zero row or a pivot too small relative to its
row scale raises \api{std::domain_error}, identifying a singular or rejected
ill-conditioned system. Use finite coefficients and a sensible finite
nonnegative pivot threshold; this is a small dense solver, not a rank-revealing
factorization or a condition-number estimator. Global interpolation uses this
reusable solver rather than duplicating it inside spline classes.

\subsection{Square, absolute comparison, and definition sizing}

\api{square(value)} returns \api{value * value}; it is \api{constexpr} and
\api{noexcept}, with ordinary floating-point range and rounding behavior.
It does not trap an overflow. \api{approximately_equal(first, second, tolerance)}
returns whether the absolute difference is at most the supplied tolerance.
It performs no relative scaling and does not validate the tolerance. Select a
nonnegative tolerance in the compared quantities' units. For example, a fixed
absolute tolerance suitable for millimeter coordinates need not suit the same
data expressed in kilometers.

\api{nurbs_knot_count(degree, control_point_count)} returns
$C+p+1$. It is shared by both worlds and both spline storage forms. The degree
must be positive and smaller than the control-point count; invalid values
raise \api{std::invalid_argument}. An unrepresentable result raises
\api{std::overflow_error}. This helper sizes the complete knot vector; it does
not generate knot values or decide whether a candidate vector is clamped.
\par\noindent\begin{minipage}{\linewidth}
\begin{lstlisting}[language=C++]
constexpr std::size_t knots_needed = np::nurbs_knot_count(3, 5);
static_assert(knots_needed == 9);
const REAL area_factor = np::square(REAL(2.5));
const bool close = np::approximately_equal(
    REAL(1), REAL(1) + REAL(5.0000e-09), REAL(1.0000e-08));
// area_factor rounds to 6.2500e+00; close is true.
\end{lstlisting}
\end{minipage}\par

For two points in the same world, \api{distance_squared(first, second)} and
\api{distance(first, second)} are also available from their point headers.
They return ordinary squared Euclidean distance and Euclidean distance,
respectively, without contact snapping or a search. Unlike the target-entity
distance overloads, they need neither a tolerance nor numerical settings.

\begin{manualnote}[A practical verification pattern]
Preserve the parameters and world positions from detailed results when they
matter to downstream work. Re-evaluate the entity at those parameters, inspect
the residual, and repeat sampled queries at higher sampling density when
coverage is uncertain. Round only presentation output; keep the underlying
floating-point values for subsequent calculations and persistence.
\end{manualnote}
